\documentclass[10pt,a4paper]{amsart}
\usepackage[margin=19mm]{geometry}
\usepackage{amsmath,amssymb,mathtools}
\usepackage[scr=rsfs,cal=euler]{mathalfa}
\usepackage{xfrac}
\usepackage[unicode,hidelinks,bookmarks=false]{hyperref}
\newcommand{\ie}{i.e.\ }
\newcommand{\cf}{cf.\ }
\newcommand{\resp}{respectively\ }
\newcommand{\Subsection}{\subsection}
\providecommand{\nobreakdash}{\nobreak\hbox{-}}
\setlength{\emergencystretch}{2em}
\multlinegap0pt
\usepackage{bm}
\usepackage{bbm}
\usepackage{dsfont}
\usepackage{xspace}
\usepackage{cancel}
\usepackage{mathtools}
\usepackage{amsthm}
\makeatletter
\@ifundefined{thm}{\newtheorem{thm}{Thm}}{}
\@ifundefined{lem}{\newtheorem{lem}[thm]{Lemma}}{}
\makeatother
\def\nil{\textit{nil}}
\title[The Four Color Theorem with Linearly Many Reducible Configur]{The Four Color Theorem with Linearly Many Reducible Configurations and Near-Linear Time Coloring}
\author{Yuta Inoue, Ken-ichi Kawarabayashi, Atsuyuki Miyashita, Bojan Mohar, Carsten Thomassen, Mikkel Thorup}
\date{Source: arXiv:2603.24880v2 (CC BY 4.0)}
\begin{document}
\maketitle

\noindent\emph{Source and licence.} The Four Color Theorem with Linearly Many Reducible Configurations and Near-Linear Time Coloring. Version arXiv:2603.24880v2 (CC BY 4.0), distributed under the Creative Commons Attribution 4.0 International licence, as recorded in the arXiv licence metadata for this exact version and shown on the abstract page https://arxiv.org/abs/2603.24880v2. Full rights notice: licenses/arxiv-2603.24880-CC-BY-4.0.txt.\par\smallskip

\noindent\textit{Editorial scope.} Counted unit: the Lem whose source label is \texttt{lem:free-homo-conf} in the article (source0.tex lines 1926--1934, source label \texttt{lem:free-homo-conf}), delivered together with the article's own complete original proof of it (source0.tex lines 1936--2017, one \texttt{proof} environment of 82 lines). No mathematics is written, repaired or re-derived: statement and proof are the article's own text line for line. Required context retained uncounted: lem:free-homo-trian (lines 1833--1841). Every definition, display and label of those lines is retained unchanged. Omitted from this package and not redistributed: 1--1832, 1842--1925, 1935--1935, 2018--3651. Nothing inside a retained excerpt is omitted or altered: every retained body is delivered exactly as the article's lines are written, so no omission receipt of any kind accompanies this package. Citations: the retained excerpts carry no citation command at all, so no bibliography entry of the article is transcribed and no reference is redistributed. Reference adaptation: none was needed and none was made; every reference inside the delivered unit resolves inside it, so no unresolved reference or citation is produced. Printed slips kept literal: no printed word, symbol, hypothesis or number of the counted statement or of its proof was altered. Layout and class: the article's own class is \texttt{article} with packages including graphicx, fullpage, inputenc, amsmath, amsfonts, amsthm, amssymb, graphics, url, hyperref, color, enumitem, lineno, subfiles; the wrapper is the lane's pinned \texttt{amsart} editorial wrapper at 10pt on A4, which restates the theorem-like environments the retained text needs and adds the article's own macro definitions that the retained text needs (\texttt{\textbackslash nil}), with extra wrapper packages (bm, bbm, dsfont, xspace, cancel, mathtools). The delivered numbering is the unit's own. Assets: no retained span contains a figure environment, a graphics-inclusion command, a PDF-page inclusion, a table or a TikZ picture; no image file is copied into this package, the package declares no asset dependency and no image file is redistributed with it. Licence: version arXiv:2603.24880v2 of this article is distributed under the Creative Commons Attribution 4.0 International licence, as recorded in the arXiv licence metadata for this exact version and shown on the abstract page https://arxiv.org/abs/2603.24880v2. A full rights notice is delivered at licenses/arxiv-2603.24880-CC-BY-4.0.txt. Characters: every retained excerpt is pure ASCII, so the delivered engine, XeLaTeX with its default Latin Modern text font, reports no missing character.
\begin{lem}\label{lem:free-homo-trian}
Suppose that we are given a pseudo-triangulation $Z$ and a set $S$ of dart identification requests. Our construction yields
a homomorphism $\phi^*$ from $Z$ to a a minimal image $Z^*$
that is a pseudo-triangulation. Both $\phi^*$ and $Z^*$ are unique up to an isomorphism. 
Moreover, consider any other homomorphism $\phi'$ from $Z$ to a 
pseudo-triangulation $Z'$ that respects $S$. Then there
exists a homomorphism $\phi^{*\prime}$ from $Z^*$ to $Z'$ such that 
$\phi^{*\prime}\circ\phi^*=\phi'$.
\end{lem}

 \begin{lem}\label{lem:free-homo-conf}
Suppose that we are given a pseudo-configuration $(Z,\delta)$ and a set $S$ of dart identification requests. Suppose there is a homomorphism $\phi'$ from $(Z,\delta)$ to a pseudo-configuration $(Z',\delta')$ such that $\phi'$ respects $S$.
Then we can 
successfully construct a free homomorphism $\phi^*$ with image $(Z^*,\delta^*)$ respecting $S$. Up to an isomorphism,
the image $(Z^*,\delta^*)$ 
is unique, independent of  $\phi'$ and $(Z',\delta')$, and it is guaranteed that there
exists a homomorphism $\phi^{*\prime}$ from $(Z^*,\delta^*)$ to $(Z',\delta')$ so that 
$\phi^{*\prime}\circ\phi^*=\phi'$.
\end{lem}

\begin{proof}
The first part of our construction is to ignore the degrees and just apply Lemma \ref{lem:free-homo-trian} to $Z$ and $Z'$. This yield a homomorphism $\phi^*$ with  image $Z^*$ respecting $S$ and with a homomorphism $\phi^{*\prime}$ from 
$Z^*$ to $Z'$ such that
$\phi^{*\prime}\circ\phi^*=\phi'$.

Now that $S$ is respected, the rest of the construction
aims to manipulate $Z^*$ to become a degree-respecting configuration 
$(Z^{**},\delta^{**})$ (which will play the role of $(Z^*,\delta^*)$ in the statement of Lemma \ref{lem:free-homo-conf}), 
or produce an error if this is not possible. However, here $(Z',\delta')$ serves as a witness that the construction is possible, hence that an error should not occur.

As we do the modifications to $Z^*$
we will maintain that $Z^*$
is a pseudo-triangulation with homomorphisms $\phi^*$ from $Z$ to $Z^*$ and $\phi^{*\prime}$ from $Z^*$ to $Z'$ such that
$\phi^{*\prime}\circ \phi^*=\phi'$.

Assuming the above homomorphisms, we claim that  $Z^*$ cannot have a loop; for since $Z'$ is in a configuration, it has no loops, but then $Z^*$ cannot have a loop
since $\phi^{*\prime}$ would map such a loop in $Z'$.

Likewise we argue that there cannot be a degree mismatch; for consider
any two vertices $u$ and $v$ with $\phi^*(u)=\phi^*(v)$. Then we must
have $\phi'(u)=
\phi^{*\prime}(\phi^*(u))=
\phi^{*\prime}(\phi^*(v))=\phi'(v)$, and the definition of homorphisms between configurations imply that 
$\delta(u)=\delta'(\phi'(u))=\delta'(\phi'(v))=\delta(v)$.


Let us now consider an inner vertex $v$ in 
$Z^*$. It has a cyclic list $u_1v,\ldots,u_dv$ of darts around $v$ with
$d=d_{Z^*}(v)$. Because $\phi^{*\prime}$ is a homomorphism, $\phi^{*\prime}(u_1v),\ldots,\phi^{*\prime}(u_dv)$ must be a cyclic walk around $\phi^{*\prime}(v)$ following successors in $Z'$. Therefore $\phi^{*\prime}(v)$ is an inner vertex in $Z'$ with $d_{Z'}(\phi^{*\prime}(v))\leq d_{Z^*}(v)$. But since $Z'$ is a configuration, we have $d_{Z'}(\phi^{*\prime}(v))=\delta'(\phi^{*\prime}(v))=\delta^*(v)$. Thus, we have a contradiction 
if we report an
error saying $d_{Z^*}(v)<\delta^*(v)$.

Now let $v$ be any vertex in 
$Z^*$ with $d_{Z^*}(v)>\delta^*(v)$. Let $uv$ and $wv$ the the incident darts from our construction that are exactly $\delta^*(v)$ apart in the successor order around $v$.
Then in $Z'$ we must get from $\phi^{*\prime}(uv)$ to $\phi^{*\prime}(wv)$
following successors $\delta^*(v)$ times around $\phi^{*\prime}(v)$. However, $\delta^*(v)=\delta'(\phi^{*\prime}(v))$, and since $Z'$ is a configuration, this implies  $\phi^{*\prime}(uv)=\phi^{*\prime}(wv)$.

We know that $\phi^{*\prime}(uv)=\phi^{*\prime}(wv)$, so as in our construction, we can apply Lemma \ref{lem:free-homo-trian} with $Z=Z^*$ and $\phi'=\phi^{*\prime}$ and an input identification request $(uv,wv)$.
The result is a homomorphic image $Z^{**}$ with homomorphisms $\phi^{**}$ from $Z^*$ to $Z^{**}$ and $\phi^{**'}$ from
$Z^{**}$ to $Z'$. We can now let
$Z^{**}$ play the role of $Z^*$,
with $\phi^{**}\circ \phi^*$ the new $\phi^*$ and $\phi^{**\prime}$ the new $\phi^{*\prime}$.

Finally, we have the last case in our construction with a boundary vertex  $v$ of $Z^*$
where $uv$ and $wv$ are the first and the last in the successor list of darts with head $v$ and the length of the list is exactly $\delta^*(v)$. The image in $Z'$ with the homomorphy $\phi^{*\prime}$ has the same successor sequence
from $\phi^{*\prime}(uv)$ to $\phi^{*\prime}(wv)$ of length $\delta^*(v)$ but since $Z'$ is a configuration and $\delta'(\phi^{*\prime}(v))=\delta'(v)$, we conclude that the list is cyclic in $Z'$ with $\phi^{*\prime}(uv)$ the successor of $\phi^{*\prime}(wv)$.

Since $Z'$ is loop free, we know it has an inner
facial triangle with 
distinct vertices $u'=\phi^{*\prime}(u)$, $v'=\phi^{*\prime}(v)$, and $w'=\phi^{*\prime}(w)$ including an
edge (consisting of two darts) between $u'$ and $w'$. 
This implies that $u$, $v$, and $w$ are distinct,
so we have a contradiction if we report a boundary error because they are not distinct.

In our construction, we  added the boundary dart  $uw$ with reversal $wu$ to $Z^*$, assigning all the relevant predecessors and successors so that the homomorphism $\phi^{*\prime}$ can be extended to map $uw$ and
$wu$ to the darts in $u'w'$ and $w'u'$. The homomorism $\phi^*$ from
$Z$ to $Z^*$ is not affected by this addition,
for in the image of $Z$,
we only changed boundary  $\nil$-pointers from $uv$
and $wv$.

We have thus proved that all our modifications
of $Z^*$ maintain the homomorphisms $\phi^*$ and $\phi^{*\prime}$. It is also clear that if we terminate, it is because we respect the degrees in $\delta^*$ so that
$(Z^*,\delta^*)$ is a pseudo-configuration.


The next question is whether we are sure to terminate. First consider the case where we add a new boundary dart $uw$ with reversal $wu$ around $v$ that thereby becomes inner.
The homomorphic image of an inner vertex cannot be a boundary vertex, so each vertex can only turn from boundary to inner once. The number of boundary darts added is therefore at most twice the number of boundary vertices in $Z$.  All other operations are dart identifications, each of which reduces the number of dart identifiers in $Z^*$. The total number of dart identifications 
is therefore bounded by  twice the number of boundary vertices plus 
the number of dart identifiers in $Z$.

Finally, we want to argue that the outcome $(Z^*,\delta^*)$ is unique up to an isomorphism. 
Suppose we have obtained two outcomes $(Z^*_1,\delta^*_1)$ and
$(Z^*_2,\delta^*_2)$. We can then let $(Z^*_2,\delta^*_2)$ play the role of $(Z',\delta')$ for $(Z^*_1,\delta^*_1)$ and vice versa. We now know that $\phi^{*\prime}_1\circ \phi^*_1=\phi^*_2$, so any vertex from $Z$ identified by $\phi^*_1$ must also be identified by $\phi^*_1$, and vice versa. We also note from
our construction that $V(Z^*_i)=\phi_i^*(V(Z))$ (this does not hold for the darts, as we may add darts). We conclude that $|V(Z^*_1)|=|V(Z^*_2)|$. This further implies that $\phi^{*\prime}_1$ must be 1-1 from $V(Z^*_1)$ to $V(Z^*_2)$, but then 
$\phi^{*\prime}_1$ is also
1-1 from $E(Z^*_1)$
to $E(Z^*_1)$, implying $|E(Z^*_1)|\leq |E(Z^*_2)|$. A symmetric argument implies
$|E(Z^*_2)|\leq |E(Z^*_1)|$, so
$|E(Z^*_1)|=|E(Z^*_2)|$, and therefore $\phi^{*\prime}_1$
is a bijective homomorphism, which is equivalent to saying it is an isomorphism.
\end{proof}

\end{document}
